%% tags: [recursion]
%% source: 2023-fa-midterm_01
\begin{prob}
    Professor Billy claims to have invented a new sorting algorithm:
    \mintinline{python}{mediansort}. The algorithm supposedly works by finding a median
    number, swapping it into the middle of the list, and recursing on the left
    half and the right half. Professor Billy's code is written below:

    \inputminted[lastline=18]{python}{code.py}

    The code is supposed to sort ``in place'' by modifying the input array.
    You may assume that \mintinline{python}{find_median} correctly returns the index
    of the median element in \mintinline{python}{arr[start:stop]}.

    Does Professor Billy's \mintinline{python}{mediansort} algorithm work? You may assume that it is called
    with \mintinline{python}{start = 0} and \mintinline{python}{stop = len(arr)}.

    \begin{choices}
        \choice Yes: \mintinline{python}{arr} will be in sorted order after the function exits.
        \choice No: it may recurse infinitely.
        \choice No: it may try to access the array at an invalid index and then raise an exception.
        \correctchoice No: it will run without error, but \mintinline{python}{arr} may not be sorted when the function exits.
    \end{choices}

    \begin{soln}

        \textbf{No: it will run without error, but \mintinline{python}{arr} may not be sorted 
        when the function exits}.


        Remember the three things to check when evaluating recursive code: 1)
        does it have a base case? 2) do the recursive subproblems get smaller
        and smaller? and 3) assuming that the recursive subcalls do what they say
        they will, does the algorithm work overall?

        In this situation, the third thing does not check out. That is: even if
        \mintinline{python}{mediansort(arr, start, middle_ix)} correctly sorts the numbers
        in \mintinline{python}{arr[start:middle_ix]} (and the other call does the same for
        the right half of the list), this does not mean that \mintinline{python}{arr} will
        be in sorted order over all. We would still need to merge the two
        sorted halves together into one sorted list.

        It might help to look at this another way. On the root call to
        \mintinline{python}{mediansort}, the median is put into its correct position in the
        sorted order. But on the recursive calls, this is not true. Consider
        the recursive call on the left half -- during this call, the median of
        the first $n/2$ numbers in the list is placed into position $n/4$
        (roughly). This isn't quite what we want -- we want to put the median
        of the \emph{smallest} $n/2$ numbers here. But the first $n/2$ numbers
        in the list aren't necessarily the \emph{smallest} $n/2$ numbers in the
        list overall.

    \end{soln}

\end{prob}
